\subsection{ARM In Education}
\label{chap:Back_ARM_Education}

Nowadays, ARM is the world's leading microcontroller used in the designing of 
embedded Systems and can be found in applications ranging from smartphones 
and IoT devices to automotive and aerospace systems. According to the ARM 
Architecture Organization, in the fourth quarter of 2020, ARM produced approximately 
6.7 billion ARM-based chips, its highest ever\cite{ArmRecord}. The importance of the ARM 
architecture in today's market justifies its study in universities, given the high likelihood 
that students will use it in the job market, but is not the only one.

The simplicity of its design, the RISC-basedinstruction set and its open-source 
nature make the use of ARM-based processors ideal for teaching embedded systems 
concepts. Its instruction set allows students to easily understand the relationship 
between high-level programming constructs, written in C or assembly 
language, and low-level hardware execution, fostering a deep understanding of 
processor behavior, register manipulation, and memory management. Additionally, 
once students master the fundamental ARM instruction set and addressing modes, 
they can quickly extrapolate this knowledge to various ARM implementations, 
accelerating the learning curve across the entire ARM ecosystem.